\documentclass[11pt]{article}
\usepackage[a4paper,margin=22mm]{geometry}
\usepackage{fontspec}
\setmainfont{DejaVu Serif}
\setsansfont{DejaVu Sans}
\usepackage{microtype}
\usepackage{xcolor}
\usepackage{hyperref}
\usepackage{enumitem}
\usepackage{parskip}
\definecolor{Accent}{HTML}{971D32}
\definecolor{Muted}{HTML}{666666}
\hypersetup{colorlinks=true,urlcolor=Accent,linkcolor=Accent}
\setlist{leftmargin=1.6em,itemsep=0.3em,topsep=0.35em}
\setcounter{tocdepth}{2}
\renewcommand{\contentsname}{Contents}
\newcommand{\sectionrule}{\par\vspace{-0.5em}\noindent\textcolor{Accent}{\rule{\linewidth}{0.6pt}}\par}
\begin{document}

\begin{center}
{\sffamily\bfseries\color{Accent} TOEFL WORD OF THE DAY\par}
\vspace{8mm}
{\fontsize{42}{48}\selectfont\bfseries sanction\par}
\vspace{2mm}
{\Large\sffamily\color{Accent} /ˈsæŋkʃən/\par}
\vspace{2mm}
{\small\sffamily Local audio phrase: sanction\par}
{\small\sffamily Audio file: audios/sanction.mp3\par}
\vspace{2mm}
{\large\sffamily\color{Muted} countable or uncountable noun; transitive verb\par}
\vspace{5mm}
{\sffamily official approval \textbullet\ official penalty \textbullet\ officially approve \textbullet\ officially punish\par}
\vspace{6mm}
{\large Sanction is unusual because it can refer either to official approval or to official punishment, so context is essential.\par}
\vspace{6mm}
\begin{minipage}{0.84\textwidth}
\itshape “The government imposed economic sanctions on the country.”
\end{minipage}\par
\vspace{10mm}
\textbf{Date:} August 27th 2026\quad\textbullet\quad \textbf{Level:} B1--mid-B2 TOEFL
\vfill
Created and compiled by \textbf{Amir Kooshky}\par
\vspace{2mm}
\href{https://t.me/KooshkyTOEFL}{Telegram Channel}\quad\textbullet\quad
\href{https://www.instagram.com/kooshkytoefl}{Instagram}\quad\textbullet\quad
\href{https://t.me/KooshkyTOEFL\_pv}{Personal Telegram}
\end{center}
\newpage
\tableofcontents
\newpage

\section{Meanings and usage}
\sectionrule
\subsection{Meaning 1: Official approval or permission}
\textbf{Part of speech:} countable or uncountable noun

\textbf{IPA:} /ˈsæŋkʃən/

\textbf{Pronunciation phrase:} sanction

\textbf{Local audio file:} audios/sanction.mp3

\textbf{Register:} formal; common in legal, governmental, institutional, and political contexts

\textbf{Definition:} official approval or permission for an action, decision, or activity

\textbf{In simpler words:} official permission or approval

\textbf{Typical pattern:} with, without, receive, gain, or require + sanction

\subsubsection{Natural examples}
\begin{enumerate}
  \item The project could not proceed without the official sanction of the local authorities.
  \item The policy was introduced with the sanction of the university administration.
  \item Military action would require government sanction.
  \item The agreement received formal sanction from the governing board.
  \item The researchers could not begin the experiment without institutional sanction.
  \item The proposal eventually gained the sanction of senior officials.
  \item The organization acted without the sanction of its members.
\end{enumerate}

\subsubsection{Nuance and implication}
\textbf{Central nuance:} Sanction in this sense is not casual agreement. It usually means approval given by a person or institution with authority.

\textbf{Usual implication:} The action may not be considered officially legitimate or permitted until the appropriate authority approves it.

\textbf{Compared with prohibition:} Sanction gives official permission, while prohibition officially prevents an action.

\subsubsection{Grammar notes}
\textbf{How to use it:} Use sanction as a noun for official approval. It often appears after words such as receive, obtain, require, or without.

\textbf{What can follow it:} It is often followed by of when naming the approving authority, as in the sanction of the court.

\textbf{Common preposition:} Common patterns include with the sanction of, without sanction, and receive sanction from.

\textbf{Noun use:} Sanction can be used without a or an when referring generally to official approval, as in government sanction. A sanction is possible when referring to a specific act of approval, but this use is less common than the penalty meaning.

\subsubsection{Useful patterns}
\textbf{Pattern:} with the sanction of + authority

\textbf{Use:} Use this when an action has official approval.

\textbf{Example:} The program was launched with the sanction of the governing board.

\textbf{Pattern:} without sanction

\textbf{Use:} Use this when an action lacks official approval.

\textbf{Example:} The department changed the policy without sanction from senior administrators.

\textbf{Pattern:} receive or gain sanction

\textbf{Use:} Use this when a proposal or action obtains official approval.

\textbf{Example:} The project received government sanction after months of review.

\textbf{Pattern:} require sanction

\textbf{Use:} Use this when official approval is necessary before an action can occur.

\textbf{Example:} Any major change to the agreement requires board sanction.

\subsubsection{Common wrong usage}
\paragraph{Correction 1}
\textbf{Wrong:} The project received a sanction, so it was punished.

\textbf{Correct:} The project received official sanction, so it was approved.

\textbf{Why:} In this sense, sanction means approval. Because the word can also mean a penalty, context must make the intended meaning clear.

\paragraph{Correction 2}
\textbf{Wrong:} The plan proceeded without the sanction from the board.

\textbf{Correct:} The plan proceeded without the sanction of the board.

\textbf{Why:} When naming the approving authority after the noun sanction, of is a common formal pattern.

\subsection{Meaning 2: Official penalty or restriction}
\textbf{Part of speech:} countable noun

\textbf{IPA:} /ˈsæŋkʃən/

\textbf{Pronunciation phrase:} sanction

\textbf{Local audio file:} audios/sanction.mp3

\textbf{Register:} formal; very common in politics, international relations, law, business, and regulation

\textbf{Definition:} an official penalty or restriction used to punish a person, organization, or country for breaking a rule or behaving in an unacceptable way

\textbf{In simpler words:} official punishment

\textbf{Typical pattern:} impose, face, lift, tighten, ease, or violate + sanctions

\subsubsection{Natural examples}
\begin{enumerate}
  \item The government imposed economic sanctions on the country.
  \item The organization may face sanctions for violating environmental regulations.
  \item International sanctions severely limited the country's access to foreign markets.
  \item The regulator threatened additional sanctions if the company failed to comply.
  \item The court has several sanctions available for serious violations of its rules.
  \item Trade sanctions can reduce a country's ability to import certain goods.
  \item Officials debated whether stronger sanctions would change the government's behavior.
  \item The athlete faced disciplinary sanctions after breaking competition rules.
\end{enumerate}

\subsubsection{Nuance and implication}
\textbf{Central nuance:} A sanction is an official measure intended to punish or pressure someone into changing behavior.

\textbf{Usual implication:} Sanctions may involve fines, trade restrictions, travel restrictions, loss of privileges, or other penalties.

\textbf{Compared with incentive:} A sanction discourages behavior through punishment or restriction, while an incentive encourages behavior by offering a benefit.

\subsubsection{Grammar notes}
\textbf{How to use it:} This meaning is countable. Use a sanction for one penalty and sanctions for several or for a group of official restrictions.

\textbf{What can follow it:} Common modifiers include economic, financial, trade, international, disciplinary, and legal.

\textbf{Common preposition:} Use sanctions on or against a person, organization, or country, and sanctions for a violation.

\textbf{Noun use:} The plural sanctions is especially common when talking about measures taken against countries, governments, companies, or individuals.

\subsubsection{Useful patterns}
\textbf{Pattern:} impose sanctions on + country, organization, or person

\textbf{Use:} Use this when an authority officially introduces penalties.

\textbf{Example:} Several countries imposed sanctions on the regime.

\textbf{Pattern:} face sanctions

\textbf{Use:} Use this when a person or organization may receive official punishment.

\textbf{Example:} Companies that violate the law may face severe sanctions.

\textbf{Pattern:} economic or financial sanctions

\textbf{Use:} Use this for penalties that restrict trade, banking, money, or other economic activity.

\textbf{Example:} Economic sanctions reduced the country's access to international credit.

\textbf{Pattern:} lift sanctions

\textbf{Use:} Use this when official penalties are removed.

\textbf{Example:} The government agreed to lift some sanctions after the agreement was signed.

\textbf{Pattern:} sanctions against + person, group, or country

\textbf{Use:} Use this when identifying the target of official penalties.

\textbf{Example:} The organization called for stronger sanctions against companies involved in illegal trade.

\subsubsection{Common wrong usage}
\paragraph{Correction 1}
\textbf{Wrong:} The government put sanctions to the country.

\textbf{Correct:} The government imposed sanctions on the country.

\textbf{Why:} The standard pattern is impose sanctions on someone or something.

\paragraph{Correction 2}
\textbf{Wrong:} The company received sanction for breaking the law.

\textbf{Correct:} The company faced sanctions for breaking the law.

\textbf{Why:} When sanction means punishment, the plural sanctions is very common, and face sanctions is a natural collocation.

\paragraph{Correction 3}
\textbf{Wrong:} The country was punished by an economic sanctionings.

\textbf{Correct:} The country was subjected to economic sanctions.

\textbf{Why:} Use the noun sanctions, not sanctionings.

\subsection{Meaning 3: Officially approve or permit}
\textbf{Part of speech:} transitive verb

\textbf{IPA:} /ˈsæŋkʃən/

\textbf{Pronunciation phrase:} sanction

\textbf{Local audio file:} audios/sanction.mp3

\textbf{Register:} formal; common in legal, governmental, institutional, and political writing

\textbf{Definition:} to officially approve or give permission for an action, decision, or activity

\textbf{In simpler words:} officially approve

\textbf{Typical pattern:} sanction + action, plan, policy, agreement, use, or activity

\subsubsection{Natural examples}
\begin{enumerate}
  \item The committee sanctioned the use of additional funds for the project.
  \item The government refused to sanction the proposed military operation.
  \item The university sanctioned the construction of a new research facility.
  \item The board sanctioned the agreement after reviewing its financial consequences.
  \item The law does not sanction discrimination against employees on the basis of age.
  \item Senior officials sanctioned the release of the confidential documents.
  \item The governing body sanctioned the new competition rules.
  \item Critics argued that the policy effectively sanctioned behavior that had previously been prohibited.
\end{enumerate}

\subsubsection{Nuance and implication}
\textbf{Central nuance:} As a verb, sanction can mean give formal authority or permission for something to happen.

\textbf{Usual implication:} The approval usually comes from a person or organization with the power to authorize the action.

\textbf{Compared with forbid:} To sanction an action is to officially permit it, while to forbid it is to officially prevent it.

\subsubsection{Grammar notes}
\textbf{How to use it:} This verb normally takes the approved action, policy, plan, or activity directly as its object.

\textbf{What can follow it:} Common objects include plan, action, policy, agreement, operation, use, and construction.

\textbf{Position in a sentence:} The passive form is also common: The operation was sanctioned by the government.

\subsubsection{Useful patterns}
\textbf{Pattern:} sanction + action or plan

\textbf{Use:} Use this when an authority gives official approval.

\textbf{Example:} The board sanctioned the proposed expansion.

\textbf{Pattern:} sanction the use of + something

\textbf{Use:} Use this when an authority officially permits something to be used.

\textbf{Example:} The agency sanctioned the use of the new treatment under limited conditions.

\textbf{Pattern:} be sanctioned by + authority

\textbf{Use:} Use the passive when the action receiving approval is the main focus.

\textbf{Example:} The project was sanctioned by the national planning agency.

\textbf{Pattern:} refuse to sanction + action

\textbf{Use:} Use this when an authority declines to give official permission.

\textbf{Example:} The committee refused to sanction the proposed changes.

\subsubsection{Common wrong usage}
\paragraph{Correction 1}
\textbf{Wrong:} The government sanctioned to build the new facility.

\textbf{Correct:} The government sanctioned the construction of the new facility.

\textbf{Why:} Sanction normally takes a noun or noun phrase directly as its object, not to + verb.

\paragraph{Correction 2}
\textbf{Wrong:} The committee sanctioned for the project.

\textbf{Correct:} The committee sanctioned the project.

\textbf{Why:} When sanction means officially approve, put the approved thing directly after the verb.

\subsection{Meaning 4: Officially punish}
\textbf{Part of speech:} transitive verb

\textbf{IPA:} /ˈsæŋkʃən/

\textbf{Pronunciation phrase:} sanction

\textbf{Local audio file:} audios/sanction.mp3

\textbf{Register:} formal; common in legal, regulatory, professional, and international contexts

\textbf{Definition:} to officially punish a person, organization, or country for breaking a rule or behaving in an unacceptable way

\textbf{In simpler words:} officially punish

\textbf{Typical pattern:} sanction + person, company, organization, or country + for + violation

\subsubsection{Natural examples}
\begin{enumerate}
  \item The agency can sanction companies that repeatedly violate safety regulations.
  \item Several officials were sanctioned for their involvement in the illegal activity.
  \item Governments may sanction foreign organizations involved in prohibited trade.
  \item The professional association sanctioned the member for violating its ethical rules.
  \item Regulators can sanction firms that provide false information to consumers.
  \item The athlete was sanctioned after testing positive for a prohibited substance.
  \item The international body threatened to sanction members that failed to follow the agreement.
\end{enumerate}

\subsubsection{Nuance and implication}
\textbf{Central nuance:} In this sense, sanction means use official authority to punish or restrict someone.

\textbf{Usual implication:} The punishment usually follows a violation of a law, rule, agreement, or professional standard.

\textbf{Compared with authorize:} To authorize someone is to give permission or authority. To sanction someone in this punishment sense is to penalize them.

\subsubsection{Grammar notes}
\textbf{How to use it:} Put the person, organization, or country receiving the punishment directly after sanction.

\textbf{What can follow it:} Use for to state the reason for the punishment.

\textbf{Common preposition:} A common pattern is sanction someone for violating a rule.

\textbf{Position in a sentence:} The passive is very common: The company was sanctioned for breaking environmental regulations.

\subsubsection{Useful patterns}
\textbf{Pattern:} sanction + person or organization

\textbf{Use:} Use this when an authority officially punishes someone.

\textbf{Example:} The regulator sanctioned three firms involved in the scheme.

\textbf{Pattern:} sanction + someone + for + wrongdoing

\textbf{Use:} Use for to explain the reason for the punishment.

\textbf{Example:} The lawyer was sanctioned for violating professional rules.

\textbf{Pattern:} be sanctioned for + wrongdoing

\textbf{Use:} Use the passive when the punished person or organization is the focus.

\textbf{Example:} The company was sanctioned for repeatedly ignoring safety requirements.

\subsubsection{Common wrong usage}
\paragraph{Correction 1}
\textbf{Wrong:} The company was sanctioned because the government approved its project.

\textbf{Correct:} The company's project was sanctioned because the government approved it.

\textbf{Why:} Sanctioning a project usually means approving it. Sanctioning a company usually means punishing it. The object and context determine the meaning.

\paragraph{Correction 2}
\textbf{Wrong:} The regulator sanctioned to the company for the violation.

\textbf{Correct:} The regulator sanctioned the company for the violation.

\textbf{Why:} Put the person or organization being punished directly after sanction.

\section{Word family}
\sectionrule
\subsection{sanctioned}
\textbf{Part of speech:} adjective

\textbf{IPA:} /ˈsæŋkʃənd/

\textbf{Definition:} officially approved or permitted

\textbf{Relationship to the headword:} This adjective commonly describes something that has received official approval.

\textbf{Grammar pattern:} officially or government-sanctioned + activity, event, program, or action

\textbf{Usage note:} Context matters because sanctioned can also mean subjected to official penalties when it describes a person, company, or country.

\subsubsection{Examples}
\begin{enumerate}
  \item Only officially sanctioned activities are allowed in the protected area.
  \item The organization provided a list of sanctioned events.
\end{enumerate}

\section{Common collocations}
\sectionrule
\subsection{official sanction}
\textbf{Applies to:} Official approval or permission

\textbf{Meaning:} formal approval from an authority

\textbf{Example:} The project received official sanction from the ministry.

\subsection{government sanction}
\textbf{Applies to:} Official approval or permission

\textbf{Meaning:} formal approval from the government

\textbf{Example:} The operation could not proceed without government sanction.

\subsection{receive sanction}
\textbf{Applies to:} Official approval or permission

\textbf{Meaning:} obtain official approval

\textbf{Example:} The proposal received sanction from the governing board.

\subsection{economic sanctions}
\textbf{Applies to:} Official penalty or restriction

\textbf{Meaning:} official economic penalties or restrictions

\textbf{Example:} Economic sanctions restricted the country's access to international markets.

\subsection{impose sanctions}
\textbf{Applies to:} Official penalty or restriction

\textbf{Meaning:} officially introduce penalties

\textbf{Pattern:} impose sanctions on + person, organization, or country

\textbf{Example:} Several governments imposed sanctions on the regime.

\subsection{face sanctions}
\textbf{Applies to:} Official penalty or restriction

\textbf{Meaning:} risk receiving official punishment

\textbf{Example:} Companies that violate the regulations may face sanctions.

\subsection{lift sanctions}
\textbf{Applies to:} Official penalty or restriction

\textbf{Meaning:} officially remove penalties or restrictions

\textbf{Example:} The countries agreed to lift some sanctions after negotiations.

\subsection{tighten sanctions}
\textbf{Applies to:} Official penalty or restriction

\textbf{Meaning:} make existing penalties or restrictions stronger

\textbf{Example:} The government threatened to tighten sanctions if the violations continued.

\subsection{sanction an action}
\textbf{Applies to:} Officially approve or permit

\textbf{Meaning:} give official approval to an action

\textbf{Example:} The committee refused to sanction the proposed action.

\subsection{sanction the use of}
\textbf{Applies to:} Officially approve or permit

\textbf{Meaning:} officially allow something to be used

\textbf{Example:} The agency sanctioned the use of the new procedure.

\subsection{sanction a company}
\textbf{Applies to:} Officially punish

\textbf{Meaning:} officially punish a company

\textbf{Example:} The regulator sanctioned the company for repeated safety violations.

\subsection{be sanctioned for}
\textbf{Applies to:} Officially punish

\textbf{Meaning:} receive official punishment because of wrongdoing

\textbf{Pattern:} be sanctioned for + violation or wrongdoing

\textbf{Example:} The organization was sanctioned for breaking competition rules.

\section{Synonyms and antonyms}
\sectionrule
\subsection{Close synonyms}
\subsubsection{approval}
\textbf{Part of speech:} uncountable noun

\textbf{Applies to:} Official approval or permission

\textbf{Shared meaning:} Both can mean official permission.

\textbf{Difference:} Approval is the general word for agreeing with or permitting something. Sanction is more formal and often suggests approval from an authority.

\textbf{Example:} The project cannot begin without government approval.

\subsubsection{penalty}
\textbf{Part of speech:} countable noun

\textbf{Applies to:} Official penalty or restriction

\textbf{Shared meaning:} Both can refer to official punishment.

\textbf{Difference:} Penalty is a broad word for a punishment. Sanction often refers specifically to an official measure imposed by a government, court, regulator, institution, or other authority.

\textbf{Example:} The company received a substantial financial penalty.

\subsubsection{authorize}
\textbf{Part of speech:} transitive verb

\textbf{Applies to:} Officially approve or permit

\textbf{Shared meaning:} Both can mean officially approve an action.

\textbf{Difference:} Authorize means officially give permission or authority. Sanction is similar but is more formal and can also carry the opposite punishment meaning in other contexts.

\textbf{Example:} The director authorized the release of additional funds.

\subsubsection{penalize}
\textbf{Part of speech:} transitive verb

\textbf{Applies to:} Officially punish

\textbf{Shared meaning:} Both can mean officially punish.

\textbf{Difference:} Penalize simply means punish someone for breaking a rule or performing poorly. Sanction is more formal and usually refers to punishment by an authority or institution.

\textbf{Example:} The regulator penalized the company for failing to disclose important information.

\subsection{Direct or useful antonyms}
\subsubsection{prohibition}
\textbf{Part of speech:} countable or uncountable noun

\textbf{Applies to:} Official approval or permission

\textbf{Opposition:} Sanction in the approval sense gives official permission, while prohibition prevents or forbids an action.

\textbf{Example:} The law introduced a prohibition on the sale of the substance.

\subsubsection{forbid}
\textbf{Part of speech:} transitive verb

\textbf{Applies to:} Officially approve or permit

\textbf{Opposition:} To sanction an action is to officially permit it, while to forbid it is to officially prevent it.

\textbf{Example:} The regulations forbid construction in the protected area.

\section{Words learners may confuse}
\sectionrule
\subsection{sanction}
\textbf{Part of speech:} noun and transitive verb

\textbf{Why learners confuse them:} Sanction is a contronym: two common meanings of the same word can point in opposite directions.

\textbf{Key difference:} The same word can mean either approval or punishment. Sanction a project usually means approve it, while sanction a company or country often means punish it.

\textbf{sanction:} The board sanctioned the construction project.

\textbf{sanction:} The regulator sanctioned the company for violating safety rules.

\subsection{authorize}
\textbf{Part of speech:} transitive verb

\textbf{Why learners confuse them:} The words overlap when sanction is used in its approval sense.

\textbf{Key difference:} Authorize only means officially permit or give authority. Sanction can mean either approve or punish, depending on context.

\textbf{sanction:} The committee sanctioned the new procedure.

\textbf{authorize:} The committee authorized the use of the new procedure.

\subsection{sanctions}
\textbf{Part of speech:} plural noun versus countable or uncountable noun

\textbf{Why learners confuse them:} Different grammatical forms strongly favor different meanings, even though they come from the same word.

\textbf{Key difference:} The plural noun sanctions very often refers to penalties, especially economic or political restrictions. The singular or uncountable noun sanction can refer to official approval.

\textbf{sanction:} The government imposed sanctions on the foreign company.

\textbf{sanctions:} The operation proceeded with government sanction.

\section{Etymology}
\sectionrule
\textbf{Origin:} Sanction comes from Latin sanctio, referring to a binding law or decree.

\textbf{Development:} Because an authoritative decree could either authorize conduct or attach penalties to violations, sanction developed both its approval meaning and its punishment meaning.

\textbf{Memory link:} Think of sanction as an official act of authority: authority can either give permission or impose punishment.

\section{Practice exercises}
\sectionrule
\subsection{Word family choice}
Choose the best member of the word family for each blank.

\begin{enumerate}
  \item Only officially \_\_\_\_\_ events may be held in the protected area.
  \item The regulator may \_\_\_\_\_ companies that repeatedly violate safety rules.
\end{enumerate}

\subsection{Correct or incorrect?}
Decide whether each sentence is correct. Correct the inaccurate ones.

\begin{enumerate}
  \item The project received government sanction and was allowed to proceed.
  \item The government imposed sanctions on the country for violating the agreement.
  \item The committee sanctioned to build the new laboratory.
  \item The company was sanctioned for repeatedly breaking environmental regulations.
  \item Sanction always means punishment.
\end{enumerate}

\subsection{Collocation practice}
Complete each sentence with a natural collocation.

\begin{enumerate}
  \item Several countries decided to impose economic \_\_\_\_\_ on the regime.
  \item The project could not begin without official \_\_\_\_\_ from the ministry.
  \item Companies that violate the regulation may \_\_\_\_\_ serious sanctions.
  \item The committee refused to \_\_\_\_\_ the proposed operation.
  \item The regulator sanctioned the firm \_\_\_\_\_ misleading consumers.
\end{enumerate}

\subsection{Complete answer key}
\subsubsection{Word family choice}
\begin{enumerate}
  \item \textbf{sanctioned.} Only officially sanctioned events may be held in the protected area. \emph{Here sanctioned means officially approved.}
  \item \textbf{sanction.} The regulator may sanction companies that repeatedly violate safety rules. \emph{Here sanction is a verb meaning officially punish.}
\end{enumerate}

\subsubsection{Correct or incorrect?}
\begin{enumerate}
  \item \textbf{Correct} \emph{Here sanction means official approval.}
  \item \textbf{Correct} \emph{Here sanctions means official penalties or restrictions.}
  \item \textbf{Incorrect}. The committee sanctioned the construction of the new laboratory. \emph{Sanction normally takes a noun or noun phrase directly as its object.}
  \item \textbf{Correct} \emph{Here sanctioned means officially punished.}
  \item \textbf{Incorrect}. Sanction can mean either official punishment or official approval. \emph{The word has two common meanings that can appear opposite, so context is essential.}
\end{enumerate}

\subsubsection{Collocation practice}
\begin{enumerate}
  \item \textbf{sanctions.} Several countries decided to impose economic sanctions on the regime. \emph{Economic sanctions are official economic penalties or restrictions.}
  \item \textbf{sanction.} The project could not begin without official sanction from the ministry. \emph{Here sanction means official approval.}
  \item \textbf{face.} Companies that violate the regulation may face serious sanctions. \emph{Face sanctions means risk receiving official punishment.}
  \item \textbf{sanction.} The committee refused to sanction the proposed operation. \emph{Here sanction means officially approve.}
  \item \textbf{for.} The regulator sanctioned the firm for misleading consumers. \emph{Use sanction someone for the wrongdoing that caused the punishment.}
\end{enumerate}

\vfill
\begin{center}
\small\color{Muted} Created and compiled by \textbf{Amir Kooshky}\\
\href{https://t.me/KooshkyTOEFL}{Telegram Channel}\quad\textbullet\quad
\href{https://www.instagram.com/kooshkytoefl}{Instagram}\quad\textbullet\quad
\href{https://t.me/KooshkyTOEFL\_pv}{Personal Telegram}
\end{center}
\end{document}
